\chapter{Supporting MI Derivations}
\label{app:derivation-details}

The first two calculations support the near-independence results in
Chapter~\ref{chap:expanded} and their interpretation in
Chapter~\ref{chap:results}. The first calculation varies population
association while keeping the margins fixed. The second retains the
quadratic sampling variation needed at independence and very small MI.
The final section explains why positive MI and positive pointwise-MI
variance are equivalent when every population cell has positive probability.

\begin{samepage}
\section{Fixed-margin behaviour near independence}
\label{sec:fixed-margin-local-derivation}

Let \(P_0=(q_{ij})\) be a joint distribution in which \(X\) and \(Y\)
are independent, with positive cells
\(q_{ij}=a_i b_j\). Consider \(P_t=P_0+tH\), where every row and column
sum of \(H\) is zero. For sufficiently small \(t\), all cells remain
positive and the margins remain \(a\) and \(b\). We expand around
\(t=0\), where \(I(P_0)=0\).
\par
\end{samepage}

For \(f(t)=I(P_t)\), the second-order Taylor formula is
\begin{equation}
 f(t)=f(0)+f'(0)t+\frac{f''(0)}{2!}t^2+O(t^3).
\end{equation}
Here \(O(t^3)\) collects the omitted terms whose magnitude is bounded by
a constant times \(|t|^3\) as \(t\) approaches zero. These terms become
smaller than the retained quadratic term. The same notation below describes
the size of each omitted remainder near zero.

The fixed margins make the pointwise MI in cell \((i,j)\)
\begin{equation}
 \pmi_{P_t}(i,j)
 =\log\left(\frac{q_{ij}+tH_{ij}}{q_{ij}}\right)
 =\log\left(1+\frac{tH_{ij}}{q_{ij}}\right)
 =\frac{tH_{ij}}{q_{ij}}-
  \frac{t^2H_{ij}^2}{2q_{ij}^2}+O(t^3),
\end{equation}
using \(\log(1+z)=z-z^2/2+O(z^3)\) with
\(z=tH_{ij}/q_{ij}\). Multiplying by the cell probability gives
\begin{align*}
 (q_{ij}+tH_{ij})\pmi_{P_t}(i,j)
 &=(q_{ij}+tH_{ij})
   \left\{\frac{tH_{ij}}{q_{ij}}
    -\frac{t^2H_{ij}^2}{2q_{ij}^2}+O(t^3)\right\}\\*
 &=tH_{ij}-\frac{t^2H_{ij}^2}{2q_{ij}}
   +\frac{t^2H_{ij}^2}{q_{ij}}+O(t^3)\\*
 &=tH_{ij}+\frac{t^2H_{ij}^2}{2q_{ij}}+O(t^3).
\end{align*}
Multiplication by \(q_{ij}\) supplies the negative quadratic term;
multiplication by \(tH_{ij}\) supplies the positive one. Their combination
leaves half the positive term. When we sum over cells, the first-order
terms cancel because \(\sum_{i,j}H_{ij}=0\).
Thus \(f(0)=f'(0)=0\), \(f''(0)=\sum_{i,j}H_{ij}^2/q_{ij}\), and
\begin{equation}
 I(P_t)=\frac{t^2}{2}\sum_{i,j}\frac{H_{ij}^2}{q_{ij}}+O(t^3).
\end{equation}

To calculate \(V(P_t)\), first square the pointwise-MI expansion:
\begin{equation*}
 \pmi_{P_t}(i,j)^2
 =\left\{\frac{tH_{ij}}{q_{ij}}+O(t^2)\right\}^2
 =\frac{t^2H_{ij}^2}{q_{ij}^2}+O(t^3).
\end{equation*}
Multiplying by the cell probabilities and summing gives its second moment:
\begin{align*}
 \E_{P_t}\{\pmi_{P_t}(X,Y)^2\}
 &=\sum_{i,j}(q_{ij}+tH_{ij})
    \left\{\frac{t^2H_{ij}^2}{q_{ij}^2}+O(t^3)\right\}\\*
 &=t^2\sum_{i,j}\frac{H_{ij}^2}{q_{ij}}+O(t^3).
\end{align*}
MI itself begins with a quadratic term, so \(I(P_t)^2=O(t^4)\).
Subtracting it from this second moment leaves the quadratic term unchanged:
\begin{equation}
 V(P_t)=t^2\sum_{i,j}\frac{H_{ij}^2}{q_{ij}}+O(t^3).
\end{equation}

The row conditional mean also has no first-order term. Its denominator
remains \(a_i\) because the margins are fixed, and
\begin{align*}
 \E_{P_t}\{\pmi_{P_t}(X,Y)\mid X=i\}
 &=\frac{1}{a_i}\sum_j(q_{ij}+tH_{ij})
   \left\{\frac{tH_{ij}}{q_{ij}}+O(t^2)\right\}\\*
 &=\frac{t}{a_i}\sum_jH_{ij}+O(t^2)
 =O(t^2).
\end{align*}
The last equality uses the zero row sum of \(H\). The column mean has the
same cancellation, using \(\sum_iH_{ij}=0\).

In Equation~\eqref{eq:g-final}, the squared centred pointwise MI,
\(V(P_t)\), the conditional means, and \(I(P_t)\) all begin at order
\(t^2\). Only the unsquared pointwise-MI term in the bracket has a
first-order contribution. Therefore Equation
\eqref{eq:g-final} supplies the first-order term:
\begin{equation}
 g_{P_t}(i,j)=\frac{2tH_{ij}}{q_{ij}}+O(t^2).
\end{equation}
Since \(g_{P_t}\) has mean zero, its variance is its second moment:
\begin{align}
 \tau^2(P_t)
 &=\sum_{i,j}(q_{ij}+tH_{ij})
   \left\{\frac{2tH_{ij}}{q_{ij}}+O(t^2)\right\}^2\notag\\*
 &=4t^2\sum_{i,j}\frac{H_{ij}^2}{q_{ij}}+O(t^3).
\end{align}
Comparing the leading terms of \(I(P_t)\), \(V(P_t)\), and \(\tau^2(P_t)\)
gives \(V(P_t)\approx2I(P_t)\) and
\(\tau^2(P_t)\approx4V(P_t)\) for small positive MI along this path.
At \(t=0\), the first-order sampling variance is zero;
Appendix~\ref{sec:independence-boundary} treats that case separately.

\section{The independence boundary}
\label{sec:independence-boundary}

\subsection{Quadratic sampling error at independence}

At independence, \(p_{ij}=p_{i+}p_{+j}\), so every population
pointwise-MI value is \(\pmi_P(i,j)=\log 1=0\). Consequently,
\(I(P)=V(P)=0\), and the first-order change in MI is zero. Sampling can
still produce positive estimated MI, but its leading variation comes from
second-order cell errors.

To calculate this variation, fix a strictly positive population table
\(P\) in which \(X\) and \(Y\) are independent, and let \(h=(h_{ij})\)
describe a redistribution of probability,
with \(\sum_{i,j}h_{ij}=0\). Define
\begin{equation}
 P_\varepsilon=P+\varepsilon h,
 \qquad f(\varepsilon)=I(P_\varepsilon),
\end{equation}
for \(\varepsilon\) small enough that every cell remains positive.
The expansion point is \(\varepsilon=0\), where \(f(0)=I(P)=0\).
The second-order Taylor formula gives
\begin{equation}
 f(\varepsilon)
 =f(0)+f'(0)\varepsilon+\frac{f''(0)}{2!}\varepsilon^2
  +o(\varepsilon^2).
\end{equation}
The first derivative is
\begin{equation}
 f'(0)=\sum_{i,j}\{\pmi_P(i,j)-1\}h_{ij}
       =-\sum_{i,j}h_{ij}=0.
\end{equation}
Here the pointwise-MI terms vanish at independence, and the remaining
terms cancel because the probabilities must still sum to one.

The row and column changes are \(h_{i+}=\sum_jh_{ij}\) and
\(h_{+j}=\sum_ih_{ij}\). Applying
\(\mathrm d^2(u\log u)/\mathrm du^2=1/u\) to the joint and marginal
terms in Equation~\eqref{eq:mi-probability-sums} gives
\begin{equation}
 f''(0)=\sum_{i,j}\frac{h_{ij}^2}{p_{ij}}
        -\sum_i\frac{h_{i+}^2}{p_{i+}}
        -\sum_j\frac{h_{+j}^2}{p_{+j}}.
\end{equation}
Thus the leading change is
\begin{equation}
 I(P_\varepsilon)
 =\frac{\varepsilon^2}{2}
  \left\{\sum_{i,j}\frac{h_{ij}^2}{p_{ij}}
        -\sum_i\frac{h_{i+}^2}{p_{i+}}
        -\sum_j\frac{h_{+j}^2}{p_{+j}}\right\}
  +o(\varepsilon^2).
 \label{eq:appendix-mi-second-order}
\end{equation}
The first sum measures squared cell changes, whereas the last two remove the
contributions due to changed margins, leaving the change in association.
The remainder is negligible relative to \(\varepsilon^2\) as
\(\varepsilon\) approaches zero.

For a sample from a multinomial distribution, take \(\varepsilon h=\widehat P-P\).
Since cell-probability errors typically have size proportional to \(n^{-1/2}\),
squaring them makes \(\Ihat(P)\) proportional to \(n^{-1}\) under
independence. The quadratic form above measures the part of these errors
that changes association after accounting for the margins. An unrestricted
joint table has \(rc-1\) free probabilities because its cells sum to one.
An independent table is determined by \(r-1\) row probabilities and
\(c-1\) column probabilities. The number of remaining association
directions is therefore
\begin{equation}
 d=(rc-1)-\{(r-1)+(c-1)\}=(r-1)(c-1).
 \label{eq:independence-association-dimension}
\end{equation}
For example, start with a \(2\times2\) independence table whose four cell
probabilities are all \(1/4\). Keeping every row and column probability
equal to \(1/2\) allows tables of the form
\begin{equation*}
 \begin{pmatrix}
  1/4+\varepsilon&1/4-\varepsilon\\
  1/4-\varepsilon&1/4+\varepsilon
 \end{pmatrix},
 \qquad |\varepsilon|<1/4.
\end{equation*}
The diagonal cells gain \(\varepsilon\), while the off-diagonal cells
lose \(\varepsilon\), so every row and column total is unchanged.
Conversely, once the change in the first cell is chosen, preserving its
row total forces the opposite change in the second cell. Preserving the
first column total and the second row total then determines the other
two changes. There is therefore only one free amount of association
change, illustrating \(d=1\).

For this uniform independence table, one standardised sampling coordinate is
\begin{equation*}
 \sqrt n\left(\widehat p_{11}-\widehat p_{12}
                   -\widehat p_{21}+\widehat p_{22}\right).
\end{equation*}
The corresponding contrast of population probabilities equals
\(4\varepsilon\) for the table above, so it records the
diagonal-versus-off-diagonal change. In a sample, this contrast
is the average of observation contributions that equal \(+1\) when the
categories match and \(-1\) when they differ. Under uniform independence,
these two outcomes each have probability \(1/2\), giving mean zero and
variance one. The central limit theorem therefore gives the scaled
contrast a standard normal limiting distribution.

For general \(r\times c\) tables with fixed dimensions and positive cell
probabilities, the multinomial central limit theorem approximates the scaled
cell errors by a multivariate normal distribution. Linear combinations of these
standardised errors separate changes in the margins from changes in
association. After removing the margin changes, the remaining
\(d\) association coordinates can be chosen to have independent standard
normal limiting distributions. The quadratic form for \(2n\Ihat(P)\)
is then their sum of squares, giving the likelihood-ratio result
\begin{equation}
 2n\Ihat(P)\xrightarrow{d}\chi^2_{(r-1)(c-1)}
 \label{eq:independence-chi-squared}
\end{equation}
\citep{wilks1938,brillinger2004}. This is the usual one-table likelihood-ratio,
or \(G\), test of independence. Its chi-squared distribution comes from
quadratic cell errors, whereas the Wald test and the Expanded Welch test use the
first-order MI error. The two-sample first-order approximation therefore
does not apply unchanged at exact independence.

\subsection{Sampling variation at small positive MI}

The quadratic approximation also explains why estimated degrees of freedom
remain several units when population MI is almost zero. Keep \(r\) and
\(c\) fixed and let a positive population table approach independence with
cell departures proportional to \(n^{-1/2}\). The sampling errors, approximated
by a multivariate normal distribution, and the population departure then
contribute on the same scale. For this model based on a multinomial
distribution, also assume strictly positive limiting cell probabilities.
The association coordinates now include both sampling noise and the
scaled population departure. Write \(Z_1,\ldots,Z_d\) for independent
variables with standard normal distributions and
\(\eta_1,\ldots,\eta_d\) for the population shifts in those coordinates.
The quadratic statistic is approximated by
\begin{equation}
 \sum_{\ell=1}^d(Z_\ell+\eta_\ell)^2,
 \qquad \kappa=\sum_{\ell=1}^d\eta_\ell^2.
 \label{eq:local-association-squares}
\end{equation}
This sum has a noncentral chi-squared distribution, whose noncentrality
\(\kappa\) is the squared size of the population shift. The limiting
likelihood-ratio result under these local alternatives is given by
\citet{banerjee2005}.

To relate the population shift to MI, consider the approaching population
\(P\). Its difference from independence with the same margins is
\(p_{ij}-p_{i+}p_{+j}\), with zero row and column sums. Substituting
this difference into Equation~\eqref{eq:appendix-mi-second-order} makes
the two marginal terms vanish and gives
\begin{equation}
 2nI(P)\approx n\sum_{i,j}
       \frac{(p_{ij}-p_{i+}p_{+j})^2}{p_{i+}p_{+j}}
       \approx\kappa.
 \label{eq:local-mi-noncentrality}
\end{equation}
The weighted sum is the squared size of the scaled population association
departure, explaining the approximation \(\kappa\simeq2nI(P)\).
The MI in this local argument changes with \(n\) as the population
approaches independence. Combining the noise and population shift gives
the local approximation
\begin{equation}
 G=2n\Ihat(P)\ \dot\sim\ \chi^2_d(\kappa),
 \qquad d=(r-1)(c-1),\qquad \kappa\simeq2nI(P),
 \label{eq:local-noncentral-mi}
\end{equation}
where \(\dot\sim\) denotes an approximate distribution. At independence,
\(\kappa=0\), recovering Equation~\eqref{eq:independence-chi-squared}.

For a sampled table close to independence, apply the expansions in
Appendix~\ref{sec:fixed-margin-local-derivation} around the product of its
empirical row and column probabilities. This gives
\begin{equation}
 \Vhat(P)\approx2\Ihat(P)=\frac{G}{n},
 \qquad \widehat\tau^2(P)\approx4\Vhat(P)\approx\frac{4G}{n}.
 \label{eq:small-mi-plugin-moments}
\end{equation}
Substitution into the implemented component formula gives
\begin{equation}
 \widehat\nu_P
 \approx\frac{2n(G/n)^2}{4G/n}=\frac{G}{2}.
 \label{eq:small-mi-plugin-df}
\end{equation}
Thus sampling noise contributes to the estimated degrees of freedom even
when the population first-order approximation \(nI(P)\) is very small.

The same approximation separates that plug-in calculation from matching
the actual sampling moments. Under the approximation using a noncentral
chi-squared distribution, \(\E(G)\approx d+\kappa\) and
\(\Var(G)\approx2(d+2\kappa)\), giving
\begin{equation}
 \E\{\Vhat(P)\}\approx\frac{d+\kappa}{n},
 \qquad
 \Var\{\Vhat(P)\}\approx\frac{2(d+2\kappa)}{n^2}.
\end{equation}
Substituting these approximate sampling moments into the moment-matching
formula gives
\begin{equation}
 \nu_{P,\mathrm{mom}}
 \approx
 \frac{2\{(d+\kappa)/n\}^2}{2(d+2\kappa)/n^2}
 =\frac{(d+\kappa)^2}{d+2\kappa}.
 \label{eq:small-mi-moment-df}
\end{equation}
When \(\kappa\) is much larger than \(d\), this reduces to
\(\kappa/2\simeq nI(P)\), as in
Equation~\eqref{eq:near-independent-df}. When \(\kappa\) is small,
however, the matched degrees of freedom approach \(d\), while the mean
plug-in degrees of freedom approach \(d/2\).

For the \(4\times4\), \(n=1{,}000\), \(I(P)=0.0001\) setting,
\(d=9\) and \(\kappa\simeq0.2\). The local calculation gives mean
estimated component degrees of freedom of approximately 4.6, rather than
the population first-order value 0.1. It also predicts that the ratio of
the average sample-based variance prediction to
\(\Var\{\Vhat(P)\}\) is approximately
\begin{equation}
 \frac{4(d+\kappa)/n^2}{2(d+2\kappa)/n^2}
 =\frac{2(d+\kappa)}{d+2\kappa}\simeq1.96.
\end{equation}
This explains the scale of the estimated degrees of freedom in
Figure~\ref{fig:redesign-a4} and the roughly twofold sample-based
variance overprediction discussed alongside it.
The population formula alone omits the quadratic sampling contribution
responsible for both effects.

\section{Positive MI and pointwise-MI variance}
\label{sec:mi-variance-equivalence}

Assumption~4 in Chapter~\ref{chap:baselines} states that positive MI is
equivalent to positive pointwise-MI variance when every population cell
has positive probability. For such a table \(P=(p_{ij})\), the result is
\begin{equation}
 I(P)>0\quad\Longleftrightarrow\quad V(P)>0.
 \label{eq:positive-mi-variance-equivalence}
\end{equation}
We establish this result by showing that zero variance forces zero MI,
then considering the reverse direction. A probability-table example
illustrates the two quantities, followed by a table showing why the
positive-cell condition matters.

\subsection{Equivalence for positive probability tables}

Each cell has a pointwise-MI value that compares its probability with
the probability under independence. MI averages these values, whereas
\(V(P)\) averages their squared departures from that average:
\begin{align}
 I(P)&=\sum_{i,j}p_{ij}\pmi_P(i,j),\\*
 V(P)&=\sum_{i,j}p_{ij}\{\pmi_P(i,j)-I(P)\}^2.
 \label{eq:appendix-pointwise-variance}
\end{align}
Zero variance means that the values do not differ across possible
outcomes. In general, equal values can have a nonzero average, so
constancy alone does not show that MI is zero. The additional restriction
here comes from the cell probabilities having to sum to one.

\medskip
\thesisminiheading{From zero variance to zero MI.}

Suppose \(V(P)=0\). Every term in
Equation~\eqref{eq:appendix-pointwise-variance} is nonnegative. Because
\(p_{ij}>0\) in every cell, a nonzero squared departure in even one
cell would make the sum positive. The zero sum therefore requires
\begin{equation}
 \pmi_P(i,j)=I(P)\qquad\text{for every }(i,j).
\end{equation}
Thus every cell has the same pointwise-MI value. To determine that value,
substitute the definition of pointwise MI and exponentiate:
\begin{align}
 \log\left(\frac{p_{ij}}{p_{i+}p_{+j}}\right)&=I(P),\\*
 \frac{p_{ij}}{p_{i+}p_{+j}}&=e^{I(P)}=K,\\*
 p_{ij}&=Kp_{i+}p_{+j}.
 \label{eq:constant-pointwise-mi-ratio}
\end{align}
The number \(K\) is the same in every cell because \(I(P)\) is a
single number for the whole table. It says that every actual cell
probability is the same multiple of its probability under independence.

The products \(p_{i+}p_{+j}\) form an independence table with the same
margins as \(P\). They already have total probability one:
\begin{align}
 \sum_i\sum_jp_{i+}p_{+j}
 &=\sum_i p_{i+}\left(\sum_jp_{+j}\right)\\
 &=\sum_i p_{i+}(1)=1.
\end{align}
The first step takes \(p_{i+}\) outside the inner sum because it does
not depend on the column \(j\). The column probabilities sum to one,
and the remaining row probabilities also sum to one. Multiplying every
entry of this reference table by \(K\) therefore multiplies its total
probability by \(K\). Since the actual table must also sum to one,
Equation~\eqref{eq:constant-pointwise-mi-ratio} gives
\begin{equation}
 1=\sum_{i,j}p_{ij}
  =K\sum_{i,j}p_{i+}p_{+j}=K.
\end{equation}
Consequently \(K=1\), every cell satisfies
\(p_{ij}=p_{i+}p_{+j}\), and every pointwise-MI value is \(\log1=0\).
Their weighted average is therefore \(I(P)=0\).

\medskip
\thesisminiheading{From zero MI to zero variance.}

For the reverse direction, use the characterisation of zero MI in
Chapter~\ref{chap:background}: \(I(P)=0\) implies independence, so
\(p_{ij}=p_{i+}p_{+j}\) in every cell. Each pointwise-MI value is
then zero, giving
\begin{equation}
 V(P)=\sum_{i,j}p_{ij}\{0-0\}^2=0.
\end{equation}
The two directions establish \(I(P)=0\) if and only if \(V(P)=0\).
MI is nonnegative and variance is nonnegative, so their positive cases
are also equivalent, as stated in
Equation~\eqref{eq:positive-mi-variance-equivalence}.

\medskip
\thesisminiheading{A worked \(2\times2\) example.}

Consider the population table \(P\) and the independence table formed
from its margins:
\begin{equation}
 P=\begin{pmatrix}0.4&0.1\\0.1&0.4\end{pmatrix},
 \qquad
 P_{\mathrm{ind}}=\begin{pmatrix}0.25&0.25\\0.25&0.25\end{pmatrix}.
\end{equation}
Every row and column totals 0.5, so each independence probability is
\(0.5(0.5)=0.25\). A diagonal cell has actual-to-independence ratio
\(0.4/0.25=1.6\), giving pointwise MI \(\log1.6\simeq0.47000\).
An off-diagonal cell has ratio \(0.1/0.25=0.4\), giving pointwise MI
\(\log0.4\simeq-0.91629\). The two diagonal cells together have
probability 0.8, and the other two have probability 0.2. Hence
\begin{align}
 I(P)&=0.8\log1.6+0.2\log0.4\simeq0.19274,\\
 V(P)&=0.8\{\log1.6-I(P)\}^2
       +0.2\{\log0.4-I(P)\}^2\simeq0.30749.
\end{align}
MI is positive, and the differing pointwise-MI values give positive
variance. The diagonal probabilities are larger than independence
predicts, while the off-diagonal probabilities are smaller, preserving
total probability one. For \(P_{\mathrm{ind}}\) itself, every ratio
is one, so every pointwise-MI value is zero and both MI and the
pointwise-MI variance are zero.

\subsection{Why positive cell probabilities are required}

Now consider a population in which only the matching pairs can occur:
\begin{equation}
 P=\begin{pmatrix}0.5&0\\0&0.5\end{pmatrix}.
\end{equation}
The margins are still 0.5 and 0.5. Each possible outcome has pointwise
MI \(\log(0.5/0.25)=\log2\), so
\begin{equation}
 I(P)=\log2>0,
 \qquad
 V(P)=0.5\{\log2-\log2\}^2
      +0.5\{\log2-\log2\}^2=0.
\end{equation}
The zero-probability cells contribute no term to the variance, so the
zero-variance condition constrains only the two possible outcomes. The
ratio of two holds only for the diagonal cells, whose independence probabilities sum
to 0.5. Doubling those probabilities gives total probability one without
requiring a ratio of one in every cell. This table therefore shows that
positive MI can coexist with zero pointwise-MI variance when some
population cells are impossible. The positive-cell assumption concerns
population probabilities, so an unobserved cell in a finite sample does
not by itself violate the assumption.
